\begin{lemma}[The contradiction after Claim 4.5]
Let \(\mathcal T\) be a certified ETT satisfying
\noderef{MaximumProperty}, and let the positive zero-based slot
\(q\) represent the manuscript stage \(j-1\).  Write
\[
 T_{j-1}=\mathcal T.\mathrm{tree}(q),\qquad
 T_j=\mathcal T.\mathrm{tree}(q+1),\qquad
 D_{j-1}=\noderef{TashkinovStoredColors}(\mathcal T,q),
\]
and write \(\varphi_{j-1}\) for the displayed coloring at slot \(q\).
Let \(\mathcal L\) be the exact lower-rung ETT: its rung number is
\(q\), its terminal ladder tree, displayed coloring, stored set,
extension kind, and exact connecting-edge set are respectively
\(T_{j-1},\varphi_{j-1},D_{j-1}\), and the corresponding fields of
\(\mathcal T\) at \(q\).  Assume that \(\mathcal L\) has the maximum
property.  Retain the following two lower-rung assertions explicitly:
\begin{enumerate}
\item[(i)] every tree carrying a
  \noderef{CorrespondingExtendedTashkinovTree} certificate for
  \((\varrho_3,T_{j-1})\), whose certified replacement ETT has the
  maximum property, has a vertex set elementary under \(\varrho_3\);
\item[(vi)] every proper coloring leaving exactly \(e_0\) uncolored
  which is \((T_{j-1},D_{j-1},\varphi_{j-1})\)-stable is a
  \noderef{ModColoring} of \(\mathcal L\).
\end{enumerate}

Let \(T_j(v_n)\) be an ordered initial segment of \(T_j\), ending
exactly at \(v_n\), and let \(T'_j\) contain this same ordered seed as
an initial segment.  Assume that \(T_{j-1}\) lies entirely before
\(v_n\) in \(T_j\).  Retain the exact stage-construction alternative:
if the slot kind is PE, \(T_j(v_n)\) is obtained from \(T_{j-1}\) by a
\(\varphi_{j-1}\)-\noderef{TashkinovAugmentationSequence}; if it is SE
or RE, there are an exact connecting-edge identity \(f\) and an exact
one-edge seed \(A=T_{j-1}+f\), in the sense of
\noderef{TreeSequenceAdjoinEdge}, with the displayed connecting set
equal to \(\{f\}\), and \(T_j(v_n)\) is obtained from \(A\) by such an
augmentation sequence.  Also assume that all of \(T_j\) is obtained
from \(T_j(v_n)\) by augmentation under \(\varphi_{j-1}\).

Let \(T'_j\) be a \(\varrho_1\)-\noderef{TashkinovClosure} of the same
seed, with elementary vertex set under \(\varrho_1\).  Fix
\(u\in V(T'_j)\setminus V(T_j)\), with its exact displayed position,
and a color
\(\beta\in\overline{\varrho}_1(u)\).  Fix
\(z\in V(T_j)\), with exact position \(q_z\), such that the seed ends
no later than \(q_z\), and suppose \(v_n\) occurs strictly before
\(z\).  The choice of \(z\) has exactly the final alternative from
Claim 4.5:
\[
 q_z=\operatorname{last}(T_j),
 \quad\hbox{or}\quad
 q_z<\operatorname{last}(T_j)
       \text{ and }\beta\in\overline{\varrho}_1(z).          \tag{1}
\]

Retain the relevant Claim 4.3(b2), equation (19), and
uniquezang(iii) data in exact form:
\begin{align}
 \overline{\varrho}_1(x)&=\overline\varphi_{j-1}(x)
       &&(x\in V(T_j)\setminus\{v_n\}),                     \tag{2}\
 \varrho_1(T_j.\mathrm{edge}(a))
   &=\varphi_{j-1}(T_j.\mathrm{edge}(a))
       &&(a\text{ an exact displayed edge position}),       \tag{3}\
 \overline\varphi_{j-1}(v_n)
   &\subseteq\overline{\varrho}_1(v_n)\cup\{\gamma_{1_1}\},
                                                                    \tag{4}\
 \delta_n&\in\overline{\varrho}_1(v_n).                    \tag{5}
\end{align}
Here (4) is the explicit missing-set inclusion supplied by the
supporting-PE color chain, and (5) is the missing-color part of (19).

Finally let \(\varrho_3\) be the exact coloring selected in
\noderef{ExtensionBaseClaim45}.  Assume its complete clauses (d1)--(d3):
it is proper, leaves exactly \(e_0\) uncolored, and is
\((T_j(v_n)-v_n,D_{j-1},\varphi_{j-1})\)-stable; on the displayed
prefixes \(T_j(z)\) and \(T'_j(u)\), every missing set except possibly
those at \(u,z\), and every exact displayed edge identity, agrees with
\(\varrho_1\); moreover
\[
 \beta\in\overline{\varrho}_1(z)
     \Longrightarrow\delta_n\in\overline{\varrho}_3(z),     \tag{6}
\]
and the exact endpoint refinement is
\[
 z\in V(T'_j(u)-u)
 \Longrightarrow
 \overline{\varrho}_3(z)=\overline{\varrho}_1(z).           \tag{6z}
\]
Thus equality at \(z\) is assumed only when its displayed position in
\(T'_j\) is strictly before \(u\).  Finally,
\[
                  \gamma_{1_1}\in\overline{\varrho}_3(u).   \tag{7}
\]
The prefix agreements are expressed using
\noderef{TreeVerticesBefore} and the exact two displayed positions, so
no parallel edge identity is replaced by its endpoint pair.

Then these data are contradictory.
\end{lemma}
\begin{proof}
We first restrict (d1) to \(T_{j-1}\).  By hypothesis every vertex of
\(T_{j-1}\) occurs before \(v_n\) in \(T_j\).  Thus the missing-set
clause of \noderef{StableColoringBeforeVertex} gives
\[
 \overline{\varrho}_3(x)=\overline\varphi_{j-1}(x)
                  \qquad(x\in V(T_{j-1})).                  \tag{8}
\]
If an exact edge \(a\) is incident with \(T_{j-1}\), and its
\(\varphi_{j-1}\)-color lies in
\(\overline\varphi_{j-1}(T_{j-1})\cup D_{j-1}\), choose a witnessing
vertex of \(T_{j-1}\).  That vertex lies before \(v_n\), so the
edge clause of the same stability hypothesis fixes the individual edge
identity.  Together with (8), this is precisely
\[
 \noderef{StableColoring}
   (T_{j-1},D_{j-1},\varphi_{j-1},\varrho_3).                \tag{9}
\]
Properness, the unique-uncolored-edge assertion, (9), and lower-rung
assertion (vi) now give
\[
                       \noderef{ModColoring}(\mathcal L,\varrho_3).
                                                                    \tag{10}
\]

We next check the copied augmentations, rather than merely declaring
that they remain valid.  Consider one
\(\varphi_{j-1}\)-\noderef{TashkinovAugmentation} in the exact
stage-construction sequence ending at \(T_j(v_n)\).  Its definition
supplies its newly appended edge identity \(a\), its unchanged ordered
initial segment, a color \(c\), and an old vertex \(x\) at which \(c\)
is missing.  The displayed stage construction ends at \(v_n\), so this
old witness \(x\) occurs strictly before \(v_n\).  Apply the full
missing-set clause of (d1) at \(x\), rather than (8), to make \(c\) missing
at the same vertex under \(\varrho_3\).  The edge part of (d1) fixes the
exact identity \(a\), since its old color \(c\) is missing on the full
before-\(v_n\) set.  Hence this very one-edge
step is a \(\varrho_3\)-Tashkinov augmentation.  Induction through the
finite relation \noderef{TashkinovAugmentationSequence} copies the whole
sequence.  In the SE/RE alternative the one-edge adjoin relation is
purely structural and retains the same exact connecting edge; the
subsequent augmentation sequence is copied by the same argument.  We
therefore have the appropriate \(\varrho_3\)-eligible route from the
prescribed stage seed to \(T_j(v_n)\).                           \tag{11}

Let
\[
 U=\noderef{TreeSequencePrefix}(T'_j,q_u),
 \qquad Z=\noderef{TreeSequencePrefix}(T_j,q_z),              \tag{12}
\]
where \(q_u\) is the displayed position of \(u\).  Thus these are the
exact ordered prefixes \(T'_j(u)\) and \(T_j(z)\), including their
terminal vertices.  Unfold the closure certificate from
\noderef{TashkinovClosure}.  Its first component is a finite
\(\varrho_1\)-augmentation sequence from \(T_j(v_n)\) to \(T'_j\).
Stop it at the unique step that creates the displayed position \(q_u\).
At every retained step, the new exact edge lies strictly before \(u\)
and its missing-color witness is an old vertex.  Clause (d2) preserves
the value of that exact edge identity.  For the old missing-color witness
\(w\), the tree-sequence order puts \(w\) strictly before \(u\), so
\(w\in V(T'_j(u)-u)\).  If \(w=z\), use the conditional endpoint clause
(6z).  If \(w\ne z\), then \(w\ne u\) and the ordinary missing-set clause
of (d2) applies.  Thus in either case the complete missing set of the
actual old witness is preserved; neither the new terminal vertex \(u\) nor
the later endpoint of a one-edge step is used as its old witness.  Consequently
each retained step satisfies the two defining color clauses of
\noderef{TashkinovAugmentation} under \(\varrho_3\).  Induction gives
\[
 \noderef{TashkinovAugmentationSequence}
       (\varrho_3,T_j(v_n),U).                               \tag{13}
\]
This is the promised edge-by-edge reconstruction of the copied closure
prefix, not an assumed corresponding-ETT certificate.

We now derive the missing-set inclusion needed for the second growth.
Because \(v_n\) is strictly before \(z\) and is distinct from \(u\),
(d2) gives
\[
 \overline{\varrho}_3(v_n)=\overline{\varrho}_1(v_n).        \tag{14}
\]
Apply (4) to a color \(c\in\overline\varphi_{j-1}(v_n)\).
If \(c\in\overline{\varrho}_1(v_n)\), equation (14) makes it missing
at \(v_n\) under \(\varrho_3\).  Otherwise membership in the union in
(4) forces \(c=\gamma_{1_1}\), and (7) makes it missing at \(u\).
Thus uniquezang(iii), (d2), and (d3) give exactly
\[
 \overline\varphi_{j-1}(v_n)
   \subseteq
   \overline{\varrho}_3(v_n)\cup\overline{\varrho}_3(u).    \tag{15}
\]

Starting from \(U\), scan the exact edge positions \(a\) of \(Z\) in their
original order.  At position \(a\), let \(y_a\) be its displayed new vertex
and let \(x_a\) be the strictly earlier displayed endpoint supplied by the
attachment field of \noderef{TreeSequence}.  Compare \(y_a\) with the
vertex set of the tree sequence built so far.  Skip this position exactly
when \(y_a\) is already present.  Otherwise append the exact edge identity
at \(a\), followed by the exact new vertex \(y_a\).

We verify this construction by induction on the scanned position.  The
initial common seed, including its root and every displayed seed vertex, is
already contained in \(U\).  At position \(a\), the attachment witness
\(x_a\) occurs at an earlier displayed position.  By the induction
hypothesis, that earlier vertex was either already present when its position
was scanned or was appended then; hence \(x_a\) belongs to the current
vertex set.  In the nonskip branch \(y_a\) does not, by the branch
condition.  Thus the exact edge at \(a\) joins a current vertex to the
specified genuinely new vertex.  This exact edge identity is also new:
if it were already displayed in the current tree, its attachment there
would put both of its endpoints in the current vertex set, contrary to
\(y_a\notin V(\text{current})\).  Appending this edge and vertex therefore
preserves both injectivity fields and the attachment field of the
tree-sequence.  In the skip branch no edge is appended, so no redundant
edge with an already present displayed new vertex is introduced.  This is
an evolving-current-set scan, not a filter taken once against \(V(U)\).

In each nonskip branch, the original augmentation sequence from the seed
to \(T_j\), stopped at position \(a\), supplies a color \(c\) and an earlier
vertex \(x\) with
\[
 \varphi_{j-1}(T_j.\mathrm{edge}(a))=c,
 \qquad c\in\overline\varphi_{j-1}(x).                      \tag{16}
\]
Indeed, a position inside the common seed is necessarily skipped because
its displayed new vertex already lies in \(U\); hence every nonskipped
position is one of the one-edge steps in the stated augmentation sequence
from the seed to \(T_j\).
The new edge is before \(z\), so (3) and the edge part of (d2) preserve
this exact value under \(\varrho_3\).  If \(x\ne v_n\), then (2) and
the missing-set part of (d2) make \(c\) missing at \(x\) under
\(\varrho_3\): the witness \(x\) is strictly before the displayed new
vertex and hence before \(z\), it is not \(z\), and it is not \(u\) because
\(u\notin V(T_j)\).  The scan induction has already placed \(x\) in the
current union.  If \(x=v_n\), inclusion (15) makes \(c\) missing at either
\(v_n\) or \(u\), both of which are in the current union.  Hence (16)
always supplies an actual current missing-color witness.  The structural
and color checks are exactly the clauses of
\noderef{TashkinovAugmentation}.  This proves, one edge at a time, a
finite augmentation sequence from \(U\) to a tree \(M\) with
\[
                 V(M)=V(U)\cup V(Z).                        \tag{17}
\]
For (17), each displayed vertex of \(Z\) is either in the common root/seed,
already present at its scanned position, or appended at that position;
conversely the construction appends only displayed vertices of \(Z\).
In particular, the scan used exact edge positions and preserves distinct
parallel identities; its skip decision used the evolving new-vertex
membership, not equality of endpoint pairs.

For the precise maximum-property interface, extend \(M\) to a closure
\(C\) under \(\varrho_3\) using \noderef{TashkinovClosureExists}.
Compose its augmentation sequence with (13), the sequence just checked,
and (11).  In the PE case this makes \(C\) a closure of the exact seed
\(T_{j-1}\); in the SE/RE case it makes \(C\) a closure of the exact
one-edge seed \(T_{j-1}+f\), retaining the singleton connecting-edge
identity.  Since \(u\in V(C)\setminus V(T_j)\) while
\(V(T_{j-1})\subseteq V(T_j)\), this is a strict extension of the lower
terminal ladder tree.

Apply the universal clause of (10) to this strict eligible extension.
It yields an actual
\[
 K:\noderef{CorrespondingExtendedTashkinovTree}
          (G,k,e_0,\mathcal L,\varrho_3,C).                  \tag{18}
\]
Now apply \noderef{MaximumPropertyTransport} to \(K\), (10), and the
maximum property of \(\mathcal L\).  This proves the proposition
\[
                   \noderef{MaximumProperty}(K.\mathrm{replacementETT}),
                                                                    \tag{19}
\]
rather than treating maximum property as an informal label.  Lower-rung
assertion (i), applied to the same certificate (18) and the transported
fact (19), gives
\[
           \noderef{ColoredVertexSetElementary}
                    (G,\varrho_3,V(C)).                     \tag{20}
\]

We also apply the original maximum property of \(\mathcal T\) at the
exact positive index \(q=j-1\).  The stage equalities identify its
current tree, stored set, and reference coloring with
\(T_{j-1},D_{j-1},\varphi_{j-1}\); (9) gives its stability premise.
The two alternatives in (11) give exactly its PE or SE/RE seed premise,
including the same edge identity \(f\), and the preceding construction
gives the required closure certificate for \(C\).  The exact stage-bound
clause of \noderef{MaximumProperty} therefore yields
\[
                         |V(C)|\le |V(T_j)|.                 \tag{21}
\]

It remains only to split the alternative (1).  Suppose first that
\(q_z\) is the last position of \(T_j\).  Then \(Z=T_j\).  Equations
(17) and the closure inclusion put both \(V(T_j)\) and \(u\) in
\(V(C)\), while the stated Case 2 fact says
\(u\notin V(T_j)\).  Thus the inclusion
\(V(T_j)\subseteq V(C)\) is strict, so
\(|V(T_j)|<|V(C)|\), contradicting (21).

Otherwise \(q_z\) is not last and
\(\beta\in\overline{\varrho}_1(z)\).  Clause (6) gives
\(\delta_n\in\overline{\varrho}_3(z)\).  Since \(v_n\) is strictly
before \(z\), the two vertices are distinct; they both lie in
\(V(Z)\subseteq V(C)\).  Equations (5) and (14) also give
\(\delta_n\in\overline{\varrho}_3(v_n)\).  Hence the two distinct
vertices \(z,v_n\) of \(C\) share the missing color \(\delta_n\),
contradicting elementariness (20) by the definition of
\noderef{ColoredVertexSetElementary}.  Both alternatives are impossible,
which proves the claimed contradiction.
\end{proof}
